\documentclass[11pt,a4paper]{article}
\usepackage[utf8]{inputenc}
\usepackage[T1]{fontenc}
\usepackage{lmodern}
\usepackage[margin=24mm]{geometry}
\usepackage{amsmath,amssymb,mathtools}
\usepackage{booktabs}
\usepackage{enumitem}
\usepackage{xcolor}
\usepackage[hidelinks]{hyperref}

\newcommand{\dd}{\mathop{}\!\mathrm{d}}
\newcommand{\ii}{\mathrm{i}}
\newcommand{\id}{\mathrm{id}}
\newcommand{\R}{\mathbb{R}}
\newcommand{\Z}{\mathbb{Z}}
\newcommand{\N}{\mathbb{N}}
\newcommand{\Ord}{\operatorname{ord}}
\newcommand{\Ker}{\operatorname{Ker}}
\newcommand{\Aut}{\operatorname{Aut}}
\newcommand{\slashed}[1]{\not\!#1}
\setlength{\parindent}{0pt}
\setlength{\parskip}{6pt}
\allowdisplaybreaks

\title{Quantum Field Theory --- Lecture 2}
\author{IMAPP lecture notes}
\date{Spring 2026}
\begin{document}\maketitle

\section{Pictures of time evolution}
With a time-independent Hamiltonian, the evolution operator is
\[
U(t,t_0)=\exp\!\left[-\frac{\ii}{\hbar}H(t-t_0)\right].
\]
In the Schrödinger picture states evolve and operators are fixed. In the Heisenberg picture,
\[
|\psi\rangle_H=U^\dagger|\psi(t)\rangle_S,\qquad A_H(t)=U^\dagger A_S U,
\]
so the state is time independent and operators carry the time dependence.

For perturbation theory, split the Hamiltonian into a solvable part and an interaction,
\[
H=H_0+H_{\mathrm{int}}.
\]
The interaction picture evolves states with $H_{\mathrm{int}}$ and operators with $H_0$:
\[
|\psi(t)\rangle_I=U_0^\dagger(t,t_0)|\psi(t)\rangle_S,
\qquad A_I(t)=U_0^\dagger A_S U_0.
\]
The operator equation is
\[
\ii\hbar\frac{\dd A_I}{\dd t}=[A_I,H_0]+\ii\hbar\left(\frac{\partial A}{\partial t}\right)_I.
\]

\section{The quantum harmonic oscillator}
In dimensionless variables, $[X,P]=\ii$, and the ladder operators are
\[
a=\frac{1}{\sqrt2}(X+\ii P),\qquad a^\dagger=\frac{1}{\sqrt2}(X-\ii P),
\qquad [a,a^\dagger]=1.
\]
The number operator $N=a^\dagger a$ satisfies
\[
[N,a]=-a,\qquad [N,a^\dagger]=a^\dagger.
\]
Its eigenstates obey
\[
N|n\rangle=n|n\rangle,
\qquad a|n\rangle=\sqrt n\,|n-1\rangle,
\qquad a^\dagger|n\rangle=\sqrt{n+1}\,|n+1\rangle.
\]
The state $|0\rangle$ is the vacuum, defined by $a|0\rangle=0$. The ladder operators remove or add one quantum of excitation.

\section{From a discrete string to a field}
A discrete chain of masses and springs has a finite set of coordinates $q_i(t)$. Taking the continuum limit sends the lattice spacing $a\to0$ while keeping the mass density and elastic modulus fixed. The discrete variables become a displacement field
\[
q_i(t)\longrightarrow \phi(x,t).
\]
The action becomes an integral over spacetime,
\[
S[\phi]=\int \dd t\int \dd x\,\mathcal L(\phi,\partial_t\phi,\partial_x\phi).
\]
For a simple string, the equation of motion takes the wave-equation form
\[
\partial_t^2\phi-v^2\partial_x^2\phi=0.
\]
The important conceptual step is that a field theory has a continuous set of dynamical variables, one for each point in space.

\section*{Summary}
The interaction picture separates free evolution from interactions, while the harmonic oscillator shows how quantized excitations arise. The continuum limit replaces discrete coordinates by fields and sums by spatial integrals.
\end{document}
